% Four kinds of minimizer on one curve (Biegler Def. 2.15, p. 28).
%
%   figures/tikz/uo-minimizer-kinds.tex
%       -> media/figures/uo-minimizer-kinds.{png,svg}
%
% Lecture 14 (Unconstrained Optimality Conditions), Section II, after Definition 14.1.
%
% GEOMETRY (1 unit = 0.75 cm in x, 0.9 cm in y). The curve is built from arcs with
% horizontal tangents at every extremum, so each marked point is a true extremum of the
% drawn curve:
%   (0, 3.0) -> local min (1.6, 1.4) -> local max (3.2, 2.6) -> flat segment
%   [4.6, 6.0] at height 1.8 -> local max (7.2, 2.8) -> global min (8.8, 0.8) -> (10, 2.4).
% (1) = (8.8, 0.8): global, strict, isolated.  (2) = (1.6, 1.4): local, strict, isolated, not
% global (1.4 > 0.8).  (3) = the flat segment: every point is a local minimizer (1.8 < 2.6 and
% 1.8 < 2.8 on both sides), none is strict, none is isolated.
%
% GREYSCALE: the curve is solid black; the flat segment is a thick blue bar AND labelled (3);
% the two points are filled circles labelled (1), (2). No distinction by colour alone.
\definecolor{okabeblue}{HTML}{0072B2}

\begin{tikzpicture}[
    x=0.75cm, y=0.9cm,
    >={Latex[length=2.0mm]},
    font=\small,
    lab/.style={font=\small, inner sep=1.5pt},
  ]
  \draw[->, black!55] (-0.2,0) -- (10.4,0) node[right, font=\scriptsize] {$x$};
  \draw[->, black!55] (0,-0.2) -- (0,3.4) node[above, font=\scriptsize] {$f(x)$};
  \draw[thick]
    (0.3,3.0) to[out=-70, in=180] (1.6,1.4)
              to[out=0, in=180]   (3.2,2.6)
              to[out=0, in=180]   (4.6,1.8)
              -- (6.0,1.8)
              to[out=0, in=180]   (7.2,2.8)
              to[out=0, in=180]   (8.8,0.8)
              to[out=0, in=-110]  (10.0,2.4);
  \draw[okabeblue, line width=3.0pt] (4.6,1.8) -- (6.0,1.8);
  \fill (8.8,0.8) circle (2.2pt);
  \fill (1.6,1.4) circle (2.2pt);
  \node[lab, below] at (8.8,0.75) {(1)};
  \node[lab, below] at (1.6,1.35) {(2)};
  \node[lab, below] at (5.3,1.70) {(3)};
\end{tikzpicture}
